\chapter{Emergenz von Raum, Zeit und Metrik}

\section{Einleitung}

In der IWT sind Raum und Zeit keine fundamentalen Größen. Sie emergieren aus der Struktur und Dynamik des diskreten Informationsnetzwerks. Dieses Kapitel zeigt, wie aus der diskreten Informationsverteilung \( I_k^{(n)} \) und der Kopplungsmatrix \( K_{kl}^{(n)} \) eine effektive Geometrie entsteht.

Die zentrale Idee lautet:

\[
\text{Raum ist die effektive Metrik der diskreten Informationskopplung.}
\]

Damit wird die klassische Raumzeit der Allgemeinen Relativitätstheorie nicht verworfen, sondern als makroskopischer Grenzfall einer tieferen diskreten informationsbasierten Geometrie verstanden.

\section{Von der diskreten Informationsdynamik zur diskreten Geometrie}

Die IWT unterscheidet zwei Ebenen:

\begin{itemize}
    \item \textbf{Analoge WDBT:} Fernwirkung ohne Raummodell, rein relational,
    \item \textbf{Digitale WDBT:} Diskrete Informationsnetz, aus dem Raum emergiert.
\end{itemize}

Die analoge Theorie beschreibt direkte Wechselwirkungen, benötigt aber kein ontologisches Raumzeitkontinuum. Erst die digitale Theorie führt ein Netzwerk aus Informationsknoten ein, dessen Kopplungsstruktur eine effektive diskrete Geometrie erzeugt.

\section{Definition der diskreten Informationsmetrik}

Die diskrete Informationsmetrik entsteht aus der Sensitivität des diskreten Informations-Lagrange-Funktionals gegenüber räumlichen Änderungen der Informationsverteilung.

\subsection{Diskrete Metrik aus Funktionalableitungen}

Formal ergibt sich die diskrete Metrik aus:

\[
g_{kl}^{(n)} = \frac{\partial^2 \mathcal{F}_k}{\partial (\Delta_{kl} I^{(n)})^2},
\]

mit dem diskreten Gradienten \( \Delta_{kl} I^{(n)} = I_l^{(n)} - I_k^{(n)} \).

\subsection{Direkte Berechnung aus Kopplungsmatrix}

Alternativ kann die Metrik direkt aus der Kopplungsmatrix \( K_{kl}^{(n)} \) berechnet werden:

\[
g_{kl}^{(n)} = \frac{K_{kl}^{(n)}}{\sqrt{K_{kk}^{(n)} K_{ll}^{(n)}}}.
\]

\subsection{Interpretation der diskreten Metrik}

\begin{itemize}
    \item Große Werte \( g_{kl}^{(n)} \approx 1 \): Starke Kopplung, kleine dynamische Änderungen haben große Wirkung („steife“ Geometrie),
    \item Kleine Werte \( g_{kl}^{(n)} \approx 0 \): Schwache Kopplung, die Informationsstruktur ist „weich“,
    \item Negative Werte: Repulsive Wechselwirkung oder antikorreliertes Verhalten.
\end{itemize}

Die diskrete Metrik misst also die Steifigkeit der diskreten Informationsgeometrie.

\section{Emergenz des physikalischen Raumes aus dem diskreten Netz}

Der physikalische Raum entsteht als effektive Geometrie des diskreten Informationsnetzes.

\subsection{Diskretes Linienelement}

Das diskrete Linienelement zwischen Knoten \( k \) und \( l \) ist:

\[
ds_{kl}^{(n)} = \sqrt{g_{kl}^{(n)}} \cdot d_{kl},
\]

mit fundamentaler Distanz \( d_{kl} \) (z. B. Gitterkonstante).

\subsection{Effektive kontinuierliche Metrik}

Bei Mittelung über viele Knoten ergibt sich die effektive kontinuierliche Metrik:

\[
g_{ij}(x, t) = \lim_{\text{feine Diskretisierung}} \langle g_{kl}^{(n)} \rangle.
\]

\subsection{Diskrete Netzwerkstruktur}

In der digitalen WDBT besteht der fundamentale Zustand aus:

\begin{itemize}
    \item Knoten \( k = 1, \dots, N \) mit Informationswerten \( I_k^{(n)} \),
    \item Kopplungen \( K_{kl}^{(n)} \): gewichtete Verbindungen zwischen Knoten,
    \item Update-Regeln: rekursive Transformationen \( I_k^{(n+1)} = U\bigl(I_k^{(n)}, \{I_l^{(n)}\}\bigr) \).
\end{itemize}

Die Metrik \( g_{kl}^{(n)} \) ist die effektive Beschreibung dieser diskreten Kopplungsstruktur.

\section{Emergenz der Zeit aus diskreten Update-Schritten}

Zeit entsteht aus der Ordnung der Aktualisierungsschritte des diskreten Informationsnetzes:

\[
\{I_k^{(0)}\} \to \{I_k^{(1)}\} \to \{I_k^{(2)}\} \to \dots
\]

\subsection{Diskrete Zeit als Schrittindex}

Die fundamentale Zeit ist der diskrete Schrittindex \( n \in \mathbb{N}_0 \).

\subsection{Kontinuierliche Zeit als emergente Näherung}

Die physikalische Zeit \( t \) ist eine kontinuierliche Näherung:

\[
t \approx n \cdot T,
\]

mit fundamentalem Zeitschrift \( T \).

\subsection{Zwei diskrete Zeitstrukturen}

Die Theorie unterscheidet:

\begin{itemize}
    \item \textbf{Lokale diskrete Zeit:} Bestimmt durch Transportprozesse zwischen Nachbarknoten,
    \[
    \tau_k^{(n)} = \frac{1}{\sum_{l \in \mathcal{N}(k)} |J_{kl}^{(n)}|}.
    \]
    \item \textbf{Globale diskrete Zeit:} Bestimmt durch systemische Informationsorganisation,
    \[
    \tau_{\text{global}}^{(n)} = \frac{1}{\lambda_{\text{max}}(K^{(n)})},
    \]
    mit größtem Eigenwert \( \lambda_{\text{max}} \) der Kopplungsmatrix.
\end{itemize}

Die beobachtete Zeit ist die Überlagerung beider diskreten Strukturen.

\section{Dynamik der diskreten Informationsmetrik}

Die effektive Metrik zwischen Knotengruppen \( A \) und \( B \) entwickelt sich gemäß:

\[
g_{AB}^{(n+1)} = g_{AB}^{(n)} + T \left[ \frac{I_A^{(n)} - I_A^{(n-1)}}{T} \cdot \frac{I_B^{(n)} - I_B^{(n-1)}}{T} - \lambda \frac{\Delta^2 I_{AB}^{(n)}}{I_{AB}^{(n)}} + \mu g_{AB}^{(n)} \ln\!\bigl(1 + \gamma G \rho L^2\bigr) \right].
\]

Interpretation der Terme:

\begin{itemize}
    \item \textbf{Erster Term:} Lokale Korrelation von Informationsänderungen,
    \item \textbf{Zweiter Term:} Globale Organisation (Bohm-Struktur),
    \item \textbf{Dritter Term:} Fraktale kosmische Skalierung.
\end{itemize}

\subsection{Kontinuierlicher Grenzfall}

Für \( T \to 0 \) und feine Diskretisierung ergibt sich:

\[
\frac{d}{dt} g_{ij} = \partial_i I \partial_j I - \lambda \frac{\partial_i \partial_j I}{I} + \mu g_{ij} \ln\!\bigl(1 + \gamma G \rho L^2\bigr).
\]

Diese kontinuierliche Form ist kompakt und für analytische Berechnungen nützlich, aber die fundamentale Beschreibung bleibt die diskrete rekursive Form.

\section{Zusammenfassung}

Die diskrete Informationsgeometrie der IWT umfasst:

\begin{itemize}
    \item Eine diskrete Informationsmetrik \( g_{kl}^{(n)} \) aus der Kopplungsmatrix \( K_{kl}^{(n)} \),
    \item Einen emergenten Raum als effektive Geometrie aus der Netzwerkstruktur,
    \item Eine emergente Zeit als Schrittindex \( n \) der Updates,
    \item Eine dynamische Metrik, die aus der Informationsdynamik folgt,
    \item Einen kontinuierlichen Grenzfall, der zur bekannten Raumzeit führt.
\end{itemize}

Raum und Zeit sind in der IWT keine primitiven Größen, sondern emergente Eigenschaften der Informationsdynamik.